\startvolume{Synthesis}{Connections across the eight papers}

\lesson{SYNTHESIS / 01}{Three reconstruction contracts}{synthesis_01}
{\mathcal S_M\mathcal A_M=I,\qquad \mathcal V_g^*\mathcal V_g=I,\qquad R^{\mathsf T}R=I_{N+1}}
{The series uses three different coordinate spaces. Coefficient packets partition an index set. A normalized coherent window analyzes a function over a continuous phase plane. Casimir atlases analyze a finite multiplicity space. Each has an exact synthesis map, although its coefficients, measure, and admissible range differ. The middle panel plots coherent coefficient intensity $|A|^2$ for a log-Gamma window with shape $3/2$.}
{Split $(1,2,3,4,5,6)$ into $(1,3,5)$ and $(2,4,6)$, then interleave to recover all six entries. At $N=1$, the displayed recoupling matrix sends $(1,0)$ to $(1/2,\sqrt3/2)$; its transpose recovers the input. In the continuous case, the unit-window Moyal identity integrates squared coefficients with measure $dh\,d\eta/(2\pi)$ to recover the input norm.}
{Identify the coordinate space and its synthesis rule before transporting an identity.}
{I: exact polyphase reconstruction; VI: continuous tight frame; VII--VIII: finite reconstruction. Analytic frame identities come from the corrected shape core; the diagram does not identify these three transform ranges.}

\lesson{SYNTHESIS / 02}{Arithmetic parameters become a shape}{synthesis_02}
{\alpha=\frac{2a-1}{\beta}=2\kappa,\qquad K_+|n\rangle=\sqrt{(n+1)(n+\alpha)}\,|n+1\rangle}
{The arithmetic scale exponent fixes a positive Gamma shape. The same shape enters the log-Gamma window and the ordinary positive ladder after setting $\kappa=\alpha/2$. This is a precise parameter dictionary across IV--VII. The curves show Gamma densities in the variable $u=e^y$, so their normalization uses $du$, whereas the window density uses $dy$.}
{Take $a=2$ and $\beta=2$. Then $\alpha=3/2$, $\kappa=3/4$, and the first two raising coefficients are $\sqrt{3/2}$ and $\sqrt5$. Taking $a=3/4$, $\beta=1/2$ instead gives $\alpha=1$: the positive kernel is admissible because $a>1/2$. The original absolutely convergent signed series requires $a>1$, so that second choice belongs to the positive branch.}
{One shape links the formulas; each realization retains its own domain and measure.}
{IV--V: $\alpha=(2a-1)/\beta$; V--VI: normalized log-Gamma windows and ladder; VII: $\alpha=2\kappa$. The bridge concerns the finite-span Lie-algebra formulas, not a new global-group integration theorem.}

\lesson{SYNTHESIS / 03}{A diagonal lift removes arithmetic signs}{synthesis_03}
{E=\sum_n|v_n|^2,\qquad |v_1+v_2|^2=4\cos^2(\theta/2)\quad(v_1=1,\ v_2=e^{i\theta})}
{Time averaging the arithmetic amplitude retains its diagonal packet energy. Its squared M\"obius factor is positive, but the original signed sum uses information that this diagonal contraction removes. The illustration varies a relative phase while holding both component magnitudes fixed. The energy stays at two, while the coherent sum ranges from four to zero. A full Gram kernel may retain other phase relations; the loss here is the specified diagonal operation.}
{For $(v_1,v_2)=(1,1)$, the energy is two and the squared sum is four. For $(1,-1)$, the energy remains two and the squared sum is zero. Therefore identical positive diagonal data cannot determine the signed amplitude. This exact two-component witness explains why the positive Gamma, frame, and semigroup results do not supply the principal M\"obius cancellation estimate.}
{Record which phases an averaging or measurement operation discards.}
{IV: exact time-average identity; I--VI: signed-coherence boundary; VII--VIII: measured versus coherent recoupling. This finite phase illustration is not an arithmetic counterexample or an estimate for M\"obius sums.}

\lesson{SYNTHESIS / 04}{A binary refinement becomes a tree}{synthesis_04}
{\left(C_{n,\mathbf r}^{\mathbf s}\right)^2=\frac{n!\prod_j(s_j)_{r_j}}{(\sum_js_j)_n\prod_jr_j!},\qquad \sum_jr_j=n}
{Shape addition produces normalized two-child coefficients. Multiplying them down a refinement tree cancels every intermediate factorial and Pochhammer factor. The remaining expression depends only on leaf shapes and degrees. Nodes show both numbers: a shape is not an excitation degree. This coefficient cancellation is the shape-level analogue of composing coordinate changes across recoupling trees; the two constructions use different coefficient systems.}
{Choose shapes $(1,2,3)$, total degree three, and leaf degrees $(1,1,1)$. The left route has squared factors $9/28$ and $1/3$; the right route has $15/56$ and $2/5$. Both products are $3/28$, hence both positive amplitude coefficients equal $\sqrt{3/28}$. Enumerating all ten degree-three leaf compositions gives squared coefficients summing exactly to one. This verifies the displayed example independently of a tree drawing.}
{Cancel intermediate normalization factors before interpreting a refinement as tree independent.}
{VI: two-child refinement and Dirichlet-multinomial extension; VII--VIII: common-space tree atlases. Exact rational checks cover the stated three-child example; general cancellation is the generating-function identity in the corrected shape core.}

\lesson{SYNTHESIS / 05}{Positivity needs a specified measure}{synthesis_05}
{G_{ij}=\int x^{i+j}\,d\nu(x),\qquad \gamma q=q^{-N}\ \text{alone does not imply positive orthogonality}}
{A positive measure makes its moment matrix a Gram matrix: every quadratic form is an integral of a squared polynomial. This mechanism underlies III's positive spectral branch. A terminating scalar q-Racah expression is a separate construction. Nonzero denominators make its values defined, but do not ensure a distinct lattice, positive weights, or nonzero orthogonal norms. The right panel shows two labels landing on the same lattice point.}
{For the illustrative measure $\nu=\tfrac12\delta_{1/4}+\tfrac12\delta_{3/4}$, the first moment matrix has entries $1,1/2,5/16$ and determinant $1/16$. In VIII's exact degenerate example, take $N=1$, $q=1/2$, $\alpha=\beta=1/2$, $\gamma=4$, $\delta=1$. Then $\mu(0)=\mu(1)=3$ and $R_0=R_1=(1,1)$. Their inner product equals one for any normalized positive weights.}
{Check the lattice, weight, norms, parameter dictionary, and phases separately.}
{III: positive moment measure; VIII: scalar formulas, corrected by C11--C12. The left measure is a two-atom illustration, not the infinite arithmetic measure. The scalar counterexample is evaluated exactly with nonzero finite denominators.}

\lesson{SYNTHESIS / 06}{Match the evidence to the question}{synthesis_06}
{c_{1/2}^{\,2}=\frac4{25},\qquad s_{1/2}^{\,2}=\frac{21}{25},\qquad c_{1/2}^{\,2}+s_{1/2}^{\,2}=1}
{Three source rows ask different questions. The coefficient example compares finite reconstruction routes. The q-matrix example checks exact normalization and floating orthogonality at specified inputs. The coherence conjecture requires one constant for each positive exponent, uniformly over depth. A finite calculation preserves that distinction: it answers a bounded question without weakening the quantifiers of another claim.}
{For I.theorem.002, evaluate the six-term polynomial at $z=1/3$ directly and through two residue packets: both give $181/81$. For VIII.theorem.011 at $q=1/2$, the squared coefficients sum to one exactly; five floating-point inputs also have orthogonality residual below $10^{-14}$. The ledger marks these rows partial and queued, respectively. V.conjecture.021 remains conjectural. No Lean kernel build is performed by these calculations.}
{Attach each result to its input, arithmetic domain, source row, and verification method.}
{Preserved claim index: I.theorem.002, VIII.theorem.011, V.conjecture.021, with full IDs in the computation JSON. Status labels come from the retrieved ledger snapshot; they are not a claim that all later formalization has been searched exhaustively.}
